\expandafter\gdef\csname SNside1\endcsname{\node[anchor=north west,inner sep=0pt,text width=48mm,font=\SideFont,align=left] (SNsidebox) at ([xshift=138mm,yshift=-408.10pt]current page.north west) {\textcolor{SNink}{\textbf{课程连接 · 地址转换}}\par\vspace{1.5mm}先判断地址合法和权限，再确定驻留与映射；最后按题目模型计数访问和缺页代价。};\pgfextracty{\SNsideheight}{\pgfpointdiff{\pgfpointanchor{SNsidebox}{south}}{\pgfpointanchor{SNsidebox}{north}}}\typeout{SNHEIGHT:1:15:\the\SNsideheight}}
\expandafter\gdef\csname SNside2\endcsname{\node[anchor=north west,inner sep=0pt,text width=48mm,font=\SideFont,align=left] (SNsidebox) at ([xshift=138mm,yshift=-95.00pt]current page.north west) {\textcolor{SNink}{\textbf{课程连接 · 与组成原理按地址串联}}\par\vspace{1.5mm}组成原理研究Cache怎样定位已得到的存储地址，操作系统还要维护虚拟地址到物理页框的转换与权限。做综合访存题时先确认Cache使用虚拟还是物理字段，再排列TLB、页表、Cache和主存路径，避免混成一次查询。};\pgfextracty{\SNsideheight}{\pgfpointdiff{\pgfpointanchor{SNsidebox}{south}}{\pgfpointanchor{SNsidebox}{north}}}\typeout{SNHEIGHT:2:1:\the\SNsideheight}}
\expandafter\gdef\csname SNside3\endcsname{\node[anchor=north west,inner sep=0pt,text width=48mm,font=\SideFont,align=left] (SNsidebox) at ([xshift=138mm,yshift=-95.00pt]current page.north west) {\textcolor{SNink}{\textbf{课程连接 · 覆盖与交换}}\par\vspace{1.5mm}覆盖要求程序按预定结构安排共用内存区；交换将进程数据在内外存之间调度。};\pgfextracty{\SNsideheight}{\pgfpointdiff{\pgfpointanchor{SNsidebox}{south}}{\pgfpointanchor{SNsidebox}{north}}}\typeout{SNHEIGHT:3:16:\the\SNsideheight}}
\expandafter\gdef\csname SNside4\endcsname{\node[anchor=north west,inner sep=0pt,text width=48mm,font=\SideFont,align=left] (SNsidebox) at ([xshift=138mm,yshift=-281.76pt]current page.north west) {\textcolor{SNink}{\textbf{课程连接 · 独立条件}}\par\vspace{1.5mm}共享对象、访问权限、地址合法性和是否驻留是不同问题，不能用其中一项代替其他检查。};\pgfextracty{\SNsideheight}{\pgfpointdiff{\pgfpointanchor{SNsidebox}{south}}{\pgfpointanchor{SNsidebox}{north}}}\typeout{SNHEIGHT:4:17:\the\SNsideheight}}
\expandafter\gdef\csname SNside5\endcsname{\node[anchor=north west,inner sep=0pt,text width=48mm,font=\SideFont,align=left] (SNsidebox) at ([xshift=138mm,yshift=-95.00pt]current page.north west) {\textcolor{SNink}{\textbf{最小反例 · 空闲总量够仍不能分配}}\par\vspace{1.5mm}两个不相邻的空洞各为6MB，总空闲12MB；申请连续8MB仍失败，因为没有单个足够大的洞。只有物理相邻的空闲区才能直接合并，隔着占用区时需要搬移或改用非连续分配，不能把表中相邻行直接相加。};\pgfextracty{\SNsideheight}{\pgfpointdiff{\pgfpointanchor{SNsidebox}{south}}{\pgfpointanchor{SNsidebox}{north}}}\typeout{SNHEIGHT:5:2:\the\SNsideheight}}
\expandafter\gdef\csname SNside6\endcsname{\node[anchor=north west,inner sep=0pt,text width=48mm,font=\SideFont,align=left] (SNsidebox) at ([xshift=138mm,yshift=-505.56pt]current page.north west) {\textcolor{SNink}{\textbf{辨析 · 空闲链顺序是状态}}\par\vspace{1.5mm}最佳和最坏适应按容量排序，余区大小改变后应维护顺序；邻近适应还须保存查找位置。};\pgfextracty{\SNsideheight}{\pgfpointdiff{\pgfpointanchor{SNsidebox}{south}}{\pgfpointanchor{SNsidebox}{north}}}\typeout{SNHEIGHT:6:21:\the\SNsideheight}}
\expandafter\gdef\csname SNside7\endcsname{\node[anchor=north west,inner sep=0pt,text width=48mm,font=\SideFont,align=left] (SNsidebox) at ([xshift=138mm,yshift=-349.50pt]current page.north west) {\textcolor{SNink}{\textbf{辨析 · 空闲链顺序是状态}}\par\vspace{1.5mm}最佳和最坏适应按容量排序，余区大小改变后应维护顺序；邻近适应还须保存查找位置。};\pgfextracty{\SNsideheight}{\pgfpointdiff{\pgfpointanchor{SNsidebox}{south}}{\pgfpointanchor{SNsidebox}{north}}}\typeout{SNHEIGHT:7:21:\the\SNsideheight}}
\expandafter\gdef\csname SNside8\endcsname{\node[anchor=north west,inner sep=0pt,text width=48mm,font=\SideFont,align=left] (SNsidebox) at ([xshift=138mm,yshift=-204.42pt]current page.north west) {\textcolor{SNink}{\textbf{辨析 · 邻接按物理地址判断}}\par\vspace{1.5mm}空闲表可能按地址或容量排序。两个表项相邻并不表示物理连续；合并须检查区间端点，回收量保持守恒。};\pgfextracty{\SNsideheight}{\pgfpointdiff{\pgfpointanchor{SNsidebox}{south}}{\pgfpointanchor{SNsidebox}{north}}}\typeout{SNHEIGHT:8:12:\the\SNsideheight}}
\expandafter\gdef\csname SNside9\endcsname{\node[anchor=north west,inner sep=0pt,text width=48mm,font=\SideFont,align=left] (SNsidebox) at ([xshift=138mm,yshift=-376.19pt]current page.north west) {\textcolor{SNink}{\textbf{辨析 · 邻接按物理地址判断}}\par\vspace{1.5mm}空闲表可能按地址或容量排序。两个表项相邻并不表示物理连续；合并须检查区间端点，回收量保持守恒。};\pgfextracty{\SNsideheight}{\pgfpointdiff{\pgfpointanchor{SNsidebox}{south}}{\pgfpointanchor{SNsidebox}{north}}}\typeout{SNHEIGHT:9:12:\the\SNsideheight}}
\expandafter\gdef\csname SNside10\endcsname{\node[anchor=north west,inner sep=0pt,text width=48mm,font=\SideFont,align=left] (SNsidebox) at ([xshift=138mm,yshift=-302.27pt]current page.north west) {\textcolor{SNink}{\textbf{不变量 · 映射替换页号不改偏移}}\par\vspace{1.5mm}页大小为\(2^d\)字节且页框同样大，转换前后低d位表示同一页内位置，保持不变；高位才由虚拟页号改成物理页框号。若计算结果的低d位发生变化，应先检查是否混入了页表项地址、字段位宽或容量单位。};\pgfextracty{\SNsideheight}{\pgfpointdiff{\pgfpointanchor{SNsidebox}{south}}{\pgfpointanchor{SNsidebox}{north}}}\typeout{SNHEIGHT:10:3:\the\SNsideheight}}
\expandafter\gdef\csname SNside11\endcsname{\node[anchor=north west,inner sep=0pt,text width=48mm,font=\SideFont,align=left] (SNsidebox) at ([xshift=138mm,yshift=-95.00pt]current page.north west) {\textcolor{SNink}{\textbf{错项排除 · 多级表的收益来自稀疏}}\par\vspace{1.5mm}多级页表可不建立未使用虚拟区间对应的下级表，但每个有效映射仍要能表示页框和必要状态。若整个空间都使用，目录反而带来额外开销。因此不能说多级结构减少了每项位数，或对任何进程都一定节省总空间。};\pgfextracty{\SNsideheight}{\pgfpointdiff{\pgfpointanchor{SNsidebox}{south}}{\pgfpointanchor{SNsidebox}{north}}}\typeout{SNHEIGHT:11:4:\the\SNsideheight}}
\expandafter\gdef\csname SNside12\endcsname{\node[anchor=north west,inner sep=0pt,text width=48mm,font=\SideFont,align=left] (SNsidebox) at ([xshift=138mm,yshift=-95.00pt]current page.north west) {\textcolor{SNink}{\textbf{课程连接 · 缓存对象不同导致分支不同}}\par\vspace{1.5mm}TLB缓存地址转换，数据Cache缓存内容。TLB没找到转换时，可从主存页表补入；数据仍可能已在Cache。缺页则要检查合法页面是否不驻留。把三种事件分别列条件，才能准确计算一次访问是否需要磁盘I/O。};\pgfextracty{\SNsideheight}{\pgfpointdiff{\pgfpointanchor{SNsidebox}{south}}{\pgfpointanchor{SNsidebox}{north}}}\typeout{SNHEIGHT:12:5:\the\SNsideheight}}
\expandafter\gdef\csname SNside13\endcsname{\node[anchor=north west,inner sep=0pt,text width=48mm,font=\SideFont,align=left] (SNsidebox) at ([xshift=138mm,yshift=-95.00pt]current page.north west) {\textcolor{SNink}{\textbf{错项排除 · 多级表的收益来自稀疏}}\par\vspace{1.5mm}多级页表可不建立未使用虚拟区间对应的下级表，但每个有效映射仍要能表示页框和必要状态。若整个空间都使用，目录反而带来额外开销。因此不能说多级结构减少了每项位数，或对任何进程都一定节省总空间。};\pgfextracty{\SNsideheight}{\pgfpointdiff{\pgfpointanchor{SNsidebox}{south}}{\pgfpointanchor{SNsidebox}{north}}}\typeout{SNHEIGHT:13:4:\the\SNsideheight}}
\expandafter\gdef\csname SNside14\endcsname{\node[anchor=north west,inner sep=0pt,text width=48mm,font=\SideFont,align=left] (SNsidebox) at ([xshift=138mm,yshift=-339.17pt]current page.north west) {\textcolor{SNink}{\textbf{辨析 · 并行查询只重叠查表阶段}}\par\vspace{1.5mm}TLB与主存页表同时起查时，未命中路径已经消耗部分查表时间。命中后访问数据仍有独立费用，不能把整个访存都当成并行。};\pgfextracty{\SNsideheight}{\pgfpointdiff{\pgfpointanchor{SNsidebox}{south}}{\pgfpointanchor{SNsidebox}{north}}}\typeout{SNHEIGHT:14:13:\the\SNsideheight}}
\expandafter\gdef\csname SNside15\endcsname{\node[anchor=north west,inner sep=0pt,text width=48mm,font=\SideFont,align=left] (SNsidebox) at ([xshift=138mm,yshift=-339.17pt]current page.north west) {\textcolor{SNink}{\textbf{辨析 · 并行查询只重叠查表阶段}}\par\vspace{1.5mm}TLB与主存页表同时起查时，未命中路径已经消耗部分查表时间。命中后访问数据仍有独立费用，不能把整个访存都当成并行。};\pgfextracty{\SNsideheight}{\pgfpointdiff{\pgfpointanchor{SNsidebox}{south}}{\pgfpointanchor{SNsidebox}{north}}}\typeout{SNHEIGHT:15:13:\the\SNsideheight}}
\expandafter\gdef\csname SNside16\endcsname{\node[anchor=north west,inner sep=0pt,text width=48mm,font=\SideFont,align=left] (SNsidebox) at ([xshift=138mm,yshift=-95.00pt]current page.north west) {\textcolor{SNink}{\textbf{辨析 · 字段数与逻辑维度分开}}\par\vspace{1.5mm}段页式常编码段号、页号和偏移，但页号与偏移共同来自段内地址。教材二维地址的结论是在比较独立逻辑坐标。};\pgfextracty{\SNsideheight}{\pgfpointdiff{\pgfpointanchor{SNsidebox}{south}}{\pgfpointanchor{SNsidebox}{north}}}\typeout{SNHEIGHT:16:14:\the\SNsideheight}}
\expandafter\gdef\csname SNside17\endcsname{\node[anchor=north west,inner sep=0pt,text width=48mm,font=\SideFont,align=left] (SNsidebox) at ([xshift=138mm,yshift=-95.00pt]current page.north west) {\textcolor{SNink}{\textbf{错项排除 · 修改位不是写权限}}\par\vspace{1.5mm}可写页面刚装入且只被读取时，可以写权限为1、修改位为0；只读页被写时，应产生保护异常，不能靠换入页面解决。修改位帮助判断淘汰是否回写，权限位限制行为，两项表示的状态不是同一件事。};\pgfextracty{\SNsideheight}{\pgfpointdiff{\pgfpointanchor{SNsidebox}{south}}{\pgfpointanchor{SNsidebox}{north}}}\typeout{SNHEIGHT:17:6:\the\SNsideheight}}
\expandafter\gdef\csname SNside18\endcsname{\node[anchor=north west,inner sep=0pt,text width=48mm,font=\SideFont,align=left] (SNsidebox) at ([xshift=138mm,yshift=-95.00pt]current page.north west) {\textcolor{SNink}{\textbf{不变量 · 调页完成还要重新执行}}\par\vspace{1.5mm}合法页面不驻留时，原访存尚未成功。缺页处理把内容装入并更新映射后，通常回到原指令重试。因此费用表要区分失败查询、服务开销、重试访问；题目已包含的更新成本不能重复加，重试成本也不能漏掉。};\pgfextracty{\SNsideheight}{\pgfpointdiff{\pgfpointanchor{SNsidebox}{south}}{\pgfpointanchor{SNsidebox}{north}}}\typeout{SNHEIGHT:18:7:\the\SNsideheight}}
\expandafter\gdef\csname SNside19\endcsname{\node[anchor=north west,inner sep=0pt,text width=48mm,font=\SideFont,align=left] (SNsidebox) at ([xshift=138mm,yshift=-260.43pt]current page.north west) {\textcolor{SNink}{\textbf{错项排除 · 固定分配约束页框数量}}\par\vspace{1.5mm}本节固定分配把各进程拥有的页框数保持为既定值；若允许直接淘汰别人的页并占为己用，就改变了双方数量。题目采用可变分配时才可讨论相应全局调整，不能把分配策略与置换范围的名字任意拼接。};\pgfextracty{\SNsideheight}{\pgfpointdiff{\pgfpointanchor{SNsidebox}{south}}{\pgfpointanchor{SNsidebox}{north}}}\typeout{SNHEIGHT:19:8:\the\SNsideheight}}
\expandafter\gdef\csname SNside20\endcsname{\node[anchor=north west,inner sep=0pt,text width=48mm,font=\SideFont,align=left] (SNsidebox) at ([xshift=138mm,yshift=-95.00pt]current page.north west) {\textcolor{SNink}{\textbf{最小反例 · 命中也要更新LRU次序}}\par\vspace{1.5mm}两页框先访问A、B，再命中A，随后访问C。LRU应淘汰较久未用的B；FIFO仍可能淘汰最早装入的A。只在缺页时移动队列会遗漏命中对最近使用关系的影响，得到的就不是正确的LRU轨迹。};\pgfextracty{\SNsideheight}{\pgfpointdiff{\pgfpointanchor{SNsidebox}{south}}{\pgfpointanchor{SNsidebox}{north}}}\typeout{SNHEIGHT:20:9:\the\SNsideheight}}
\expandafter\gdef\csname SNside21\endcsname{\node[anchor=north west,inner sep=0pt,text width=48mm,font=\SideFont,align=left] (SNsidebox) at ([xshift=138mm,yshift=-95.00pt]current page.north west) {\textcolor{SNink}{\textbf{辨析 · 缺页与置换}}\par\vspace{1.5mm}空页框装入也缺页；置换须淘汰旧页。};\pgfextracty{\SNsideheight}{\pgfpointdiff{\pgfpointanchor{SNsidebox}{south}}{\pgfpointanchor{SNsidebox}{north}}}\typeout{SNHEIGHT:21:20:\the\SNsideheight}}
\expandafter\gdef\csname SNside22\endcsname{\node[anchor=north west,inner sep=0pt,text width=48mm,font=\SideFont,align=left] (SNsidebox) at ([xshift=138mm,yshift=-421.28pt]current page.north west) {\textcolor{SNink}{\textbf{辨析 · 缺页与置换}}\par\vspace{1.5mm}空页框装入也缺页；置换须淘汰旧页。};\pgfextracty{\SNsideheight}{\pgfpointdiff{\pgfpointanchor{SNsidebox}{south}}{\pgfpointanchor{SNsidebox}{north}}}\typeout{SNHEIGHT:22:20:\the\SNsideheight}}
\expandafter\gdef\csname SNside23\endcsname{\node[anchor=north west,inner sep=0pt,text width=48mm,font=\SideFont,align=left] (SNsidebox) at ([xshift=138mm,yshift=-180.56pt]current page.north west) {\textcolor{SNink}{\textbf{辨析 · 缺页与置换}}\par\vspace{1.5mm}空页框装入也缺页；置换须淘汰旧页。};\pgfextracty{\SNsideheight}{\pgfpointdiff{\pgfpointanchor{SNsidebox}{south}}{\pgfpointanchor{SNsidebox}{north}}}\typeout{SNHEIGHT:23:20:\the\SNsideheight}}
\expandafter\gdef\csname SNside24\endcsname{\node[anchor=north west,inner sep=0pt,text width=48mm,font=\SideFont,align=left] (SNsidebox) at ([xshift=138mm,yshift=-95.00pt]current page.north west) {\textcolor{SNink}{\textbf{条件 · 扫描起点决定首个候选}}\par\vspace{1.5mm}同类候选超过一个时，选择取决于题设给出的指针起点和扫描方向。};\pgfextracty{\SNsideheight}{\pgfpointdiff{\pgfpointanchor{SNsidebox}{south}}{\pgfpointanchor{SNsidebox}{north}}}\typeout{SNHEIGHT:24:23:\the\SNsideheight}}
\expandafter\gdef\csname SNside25\endcsname{\node[anchor=north west,inner sep=0pt,text width=48mm,font=\SideFont,align=left] (SNsidebox) at ([xshift=138mm,yshift=-95.00pt]current page.north west) {\textcolor{SNink}{\textbf{条件 · 扫描起点决定首个候选}}\par\vspace{1.5mm}同类候选超过一个时，选择取决于题设给出的指针起点和扫描方向。};\pgfextracty{\SNsideheight}{\pgfpointdiff{\pgfpointanchor{SNsidebox}{south}}{\pgfpointanchor{SNsidebox}{north}}}\typeout{SNHEIGHT:25:23:\the\SNsideheight}}
\expandafter\gdef\csname SNside26\endcsname{\node[anchor=north west,inner sep=0pt,text width=48mm,font=\SideFont,align=left] (SNsidebox) at ([xshift=138mm,yshift=-95.00pt]current page.north west) {\textcolor{SNink}{\textbf{课程连接 · 映射写入与磁盘持久化分开}}\par\vspace{1.5mm}程序通过映射地址修改内容，首先改变驻留页并可能置脏；何时写回文件还取决于映射属性和写回机制。文件系统维护后备关系，虚拟内存维护映射和驻留，二者协作不代表每一条赋值都会立即产生一次磁盘写入。};\pgfextracty{\SNsideheight}{\pgfpointdiff{\pgfpointanchor{SNsidebox}{south}}{\pgfpointanchor{SNsidebox}{north}}}\typeout{SNHEIGHT:26:11:\the\SNsideheight}}
\expandafter\gdef\csname SNside27\endcsname{\node[anchor=north west,inner sep=0pt,text width=48mm,font=\SideFont,align=left] (SNsidebox) at ([xshift=138mm,yshift=-95.00pt]current page.north west) {\textcolor{SNink}{\textbf{不变量 · 窗口先取引用再去重}}\par\vspace{1.5mm}工作集窗口若规定最近6次引用，应先截出准确的6次，再对页号去重。同一页出现三次仍只占集合中一项。不能先找6个不同页，也不能从当前时刻向未来补引用，否则集合不再代表题设的近期需求。};\pgfextracty{\SNsideheight}{\pgfpointdiff{\pgfpointanchor{SNsidebox}{south}}{\pgfpointanchor{SNsidebox}{north}}}\typeout{SNHEIGHT:27:10:\the\SNsideheight}}
\expandafter\gdef\csname SNside28\endcsname{\node[anchor=north west,inner sep=0pt,text width=48mm,font=\SideFont,align=left] (SNsidebox) at ([xshift=138mm,yshift=-95.00pt]current page.north west) {\textcolor{SNink}{\textbf{不变量 · 窗口先取引用再去重}}\par\vspace{1.5mm}工作集窗口若规定最近6次引用，应先截出准确的6次，再对页号去重。同一页出现三次仍只占集合中一项。不能先找6个不同页，也不能从当前时刻向未来补引用，否则集合不再代表题设的近期需求。};\pgfextracty{\SNsideheight}{\pgfpointdiff{\pgfpointanchor{SNsidebox}{south}}{\pgfpointanchor{SNsidebox}{north}}}\typeout{SNHEIGHT:28:10:\the\SNsideheight}}
\expandafter\gdef\csname SNside29\endcsname{\node[anchor=north west,inner sep=0pt,text width=48mm,font=\SideFont,align=left] (SNsidebox) at ([xshift=138mm,yshift=-95.00pt]current page.north west) {\textcolor{SNink}{\textbf{不变量 · 窗口先取引用再去重}}\par\vspace{1.5mm}工作集窗口若规定最近6次引用，应先截出准确的6次，再对页号去重。同一页出现三次仍只占集合中一项。不能先找6个不同页，也不能从当前时刻向未来补引用，否则集合不再代表题设的近期需求。};\pgfextracty{\SNsideheight}{\pgfpointdiff{\pgfpointanchor{SNsidebox}{south}}{\pgfpointanchor{SNsidebox}{north}}}\typeout{SNHEIGHT:29:10:\the\SNsideheight}}
\expandafter\gdef\csname SNside30\endcsname{\node[anchor=north west,inner sep=0pt,text width=48mm,font=\SideFont,align=left] (SNsidebox) at ([xshift=138mm,yshift=-95.00pt]current page.north west) {\textcolor{SNink}{\textbf{课程连接 · 计算顺序}}\par\vspace{1.5mm}先列地址字段、页表状态和访问分支，再计算偏移、访问次数或换页代价。};\pgfextracty{\SNsideheight}{\pgfpointdiff{\pgfpointanchor{SNsidebox}{south}}{\pgfpointanchor{SNsidebox}{north}}}\typeout{SNHEIGHT:30:18:\the\SNsideheight}}
\expandafter\gdef\csname SNside31\endcsname{\node[anchor=north west,inner sep=0pt,text width=48mm,font=\SideFont,align=left] (SNsidebox) at ([xshift=138mm,yshift=-95.00pt]current page.north west) {\textcolor{SNink}{\textbf{课程连接 · 检查清单}}\par\vspace{1.5mm}对于每次引用，说明页号与偏移、命中与缺页，以及地址转换是否需要重新执行。};\pgfextracty{\SNsideheight}{\pgfpointdiff{\pgfpointanchor{SNsidebox}{south}}{\pgfpointanchor{SNsidebox}{north}}}\typeout{SNHEIGHT:31:19:\the\SNsideheight}}
\AddToShipoutPictureFG{\begin{tikzpicture}[remember picture,overlay]\ifcsname SNside\thepage\endcsname\csname SNside\thepage\endcsname\fi\end{tikzpicture}}